\section{Related Work}
\label{sec:02_rw}

\subsection{Feature Extraction and Intermediate Processing}
\label{subsec:subsec_21_multimodal}
Transformer-based video captioning methods generally adopt one of two feature extraction paradigms: offline, where features are pre-computed and used as fixed model inputs~\cite{yangConceptAwareVideoCaptioning2023, zhongBidirectionalTransformerKnowledge2024}; and end-to-end, where the model jointly learns visual representations and caption generation from raw video inputs, e.g., SwinBERT~\cite{linSwinBERTEndtoEndTransformers2022}.

In practice, extractor outputs are typically holistic and frame-level, lacking the explicit object-level and relational structure required for fine-grained caption generation. Methods therefore introduce additional processing stages, primarily differing in the type of missing structure they target. Track4Cap~\cite{luoFramebyFrameMultiObjectTrackingGuided2025} employs a multi-object tracking module to preserve object identity across frames, and an object-relation encoder to model inter-object relationships. KG-VCN~\cite{yuanFullyExploringObject2025} models inter-object relations through a Graph Transformer alongside an external knowledge graph, while HSRA~\cite{hanActionDrivenSemanticRepresentation2025} organizes visual semantics into an object–action–event hierarchy, employing a DETR-inspired detection layer to make actions and events explicit. These modules can yield measurable gains; however, they share a trade-off: each additional component deepens the processing pipeline, increases computational overhead, and makes caption quality dependent on intermediate predictors whose errors may propagate downstream.

This trade-off becomes increasingly worth examining as modern extractors continue to improve. Pre-trained models such as CLIP~\cite{radfordLearningTransferableVisual2021} and MViTv2~\cite{liMViTv2ImprovedMultiscale2022} already produce semantically rich and well-structured representations. Given these advances, such representations may already be sufficient for direct integration, avoiding the redundant transformation and noise propagation that extensive intermediate processing may introduce.


\subsection{Semantic Knowledge Augmentation for Video Captioning}
\label{subsec:subsec_22_semantic}
Beyond structural processing, another challenge is that visual features alone often provide limited high-level semantic information and weak coverage of rare or long-tailed vocabulary. Prior work addresses this limitation through various forms of external semantic knowledge injection.

Some methods retrieve static knowledge from external knowledge bases. TextKG~\cite{guTextKnowledgeGraph2023} augments captioning with commonsense facts from ConceptNet via a two-stream framework, while BTKG~\cite{zhongBidirectionalTransformerKnowledge2024} predicts inter-object relationships using TransE and integrates the resulting relational structure into the encoder. Although these methods enrich semantic coverage, the retrieved knowledge is often only loosely conditioned on the specific video content, making relevance control a persistent challenge.

Other methods pursue scene-adaptive knowledge construction to address this relevance problem. MK-VC~\cite{sunSceneAdaptiveDynamic2026} combines static commonsense with dynamically retrieved video-related knowledge, and EMKG~\cite{sunGeneralizedVideoCaptioning2026} constructs a ConceptVision Knowledge Graph that couples detected object features with commonsense relations. While these methods improve relevance over static retrieval, they do so through more elaborate processing pipelines requiring multiple additional processing stages. Notably, MK-VC and EMKG explicitly introduce adaptive suppression mechanisms partly intended to filter noisy or irrelevant knowledge, suggesting that pipeline complexity has itself become a source of the problem it was designed to solve.

A recent study by Zhang et al.~\cite{zhangPretrainedImageTextModels2025} demonstrates that an image captioning model built on \mbox{BLIP-2~\cite{liBLIP2BootstrappingLanguageImage2023}}, post-trained on only a small set of video-text pairs, can rival specialized video captioning systems without relying on complex video-specific architectures. This finding suggests that the semantic knowledge embedded in vision-language models (VLMs) transfers effectively to the video domain, offering semantically grounded representations without the structural overhead of explicit knowledge graph pipelines.


\subsection{Decoding Strategies and Bidirectional Context Modeling}
\label{subsec:subsec_23_decoding}
Beyond feature preparation, another line of research focuses on the decoder itself. The standard Transformer decoder generates captions autoregressively from left to right (L2R), a formulation that carries two related weaknesses: the model may over-rely on language priors, and it structurally cannot condition any token on future linguistic context. Existing approaches primarily address the first weakness through various decoder-side strategies. UHCL~\cite{sunUnifiedHierarchicalContrastive2026} employs triamese decoders with hierarchical contrastive learning to reduce generic and low-distinctiveness outputs. ASGNet~\cite{liuAdaptiveSemanticGuidance2025} introduces an adaptive control decoder that dynamically rebalances visual and textual contributions at each decoding step. QPDC~\cite{chenAskFocusMore2026} uses a question-prompt mechanism to steer generation toward salient visual content. While these methods improve visual grounding, they still operate strictly from left to right and therefore cannot access future linguistic context.

Bidirectional decoding addresses this structural limitation directly. BiTransformer~\cite{zhongBiTransformerAugmentingSemantic2022} introduces a backward decoder that generates captions in a right-to-left direction, providing reverse contextual information to guide the final L2R generation. BTKG~\cite{zhongBidirectionalTransformerKnowledge2024} extends this framework by assigning separate encoders to the forward and backward decoders to accommodate different modal features. Both BiTransformer and BTKG address information leakage during training through a pseudo reverse caption strategy (see Section~\ref{subsec:subsec_46_optim}). However, both approaches still rely on encoder modules to refine extracted features, inheriting the same processing-stage trade-offs discussed in previous sections.

Our BiDecT builds upon the bidirectional decoding paradigm and pseudo reverse caption strategy of BiTransformer and BTKG, but eliminates the intermediate encoder entirely, projecting multimodal representations directly into the decoder. Additionally, BiDecT incorporates a semantic modality derived from a pre-trained vision-language model, providing semantic grounding without explicit knowledge graph construction.
